% Auto-generated from section files. Do not use \input in this template.

\documentclass[12pt,letterpaper]{article}

\usepackage[margin=1in]{geometry}
\usepackage{setspace}
\usepackage{lineno}
\usepackage{graphicx}
\usepackage{amsmath,amssymb,amsfonts}
\usepackage{amsthm}
\usepackage{booktabs}
\usepackage{siunitx}
\usepackage{hyperref}
\usepackage{natbib}

\graphicspath{{../figures/}{./}}

\theoremstyle{plain}
\newtheorem{lemma}{Lemma}

\title{A Weighted Fitting Approach for Diameter Distributions from Horizontal Point Sampling}
\author{Gregory E. Paradis\\Department of Forest Resources Management, University of British Columbia, Vancouver, BC, Canada\\\texttt{gregory.paradis@ubc.ca}}
\date{}

\begin{document}

\doublespacing

\begin{titlepage}
  \begin{center}
    {\Large\bfseries A Weighted Fitting Approach for Diameter Distributions from Horizontal Point Sampling\par}
    \vspace{1.5em}
    Gregory E. Paradis\par
    Department of Forest Resources Management\par
    University of British Columbia\par
    Vancouver, BC, Canada\par
    \texttt{gregory.paradis@ubc.ca}\par
  \end{center}
\end{titlepage}

\linenumbers

\section*{Abstract}
Horizontal point sampling (HPS) produces size-biased tallies that cannot be fit
directly with standard probability distributions without distorting diameter
distribution estimates. Previous work resolves this by deriving bespoke
size-biased probability density functions (PDFs) for each assumed distribution.
We revisit the problem and formalise a weighted non-linear least squares
approach that fits standard-form PDFs to expanded HPS stand tables while
preserving the statistical equivalence with the size-biased formulation. The
new pipeline leverages contemporary open-source software, is fully
reproducible, and includes accompanying code that regenerates all figures and
tables. Computational experiments on permanent sample plot data
show that the weighted method matches the reference approach within numerical
tolerance across Weibull and Gamma distributions. The manuscript and companion
software provide a practical workflow for practitioners who require stable,
transparent, and replicable HPS diameter distribution fitting.

\noindent\textbf{Keywords:} horizontal point sampling; diameter distribution modelling; weighted nonlinear least squares; forest biometrics; reproducible research

\newpage

\section{Introduction}

Horizontal point sampling (HPS), also referred to as angle count or Bitterlich
sampling \citep{bitterlich1947}, is widely used to quantify stand structure in
managed forests. Diameter distributions derived from HPS tallies underpin yield
modelling, inventory projection, and ecological assessments. A long-standing
challenge is that HPS tallies are intrinsically \emph{size-biased}: the
inclusion probability of a stem is proportional to its basal area, which varies
with the square of diameter. When standard probability density functions (PDFs)
are fit directly to expanded HPS stand tables, small diameter classes exert
excessive influence and the fitted distribution is biased \citep{van1986,
gove2000, ducey2015}. The conventional remedy is to derive size-biased forms of
candidate PDFs and fit them to the unexpanded tallies.

Deriving size-biased PDFs requires specialised algebra and bespoke software
implementations. Even when available, the resulting optimisation problems
often impose complex parameter bounds that reduce numerical stability. As a
result, this complexity may tempt practitioners to bypass appropriate size-bias
corrections, potentially leading to over-fitted distributions and poor downstream
predictions.

We present a simplified method for fitting diameter distributions and show
it is equivalent to the size-biased estimator under standard assumptions. This is an applied biometrics
workflow aimed at routine inventory and monitoring pipelines rather than a
theoretical development. The workflow fits standard PDFs to expanded HPS
stand tables with embedded size-bias weights, validates the approach on
permanent sample plot (PSP) meta-plots, and ships a reproducible pipeline
that regenerates figures and tables.

The remainder of the manuscript introduces notation and the reference method,
derives the weighted estimator, reports numerical comparisons, and discusses
implications for inventory analysis pipelines.

\section{Materials and methods}
\label{sec:methods}

\subsection{Notation}
Let $X$ denote the diameter at breast height (DBH) of a stem measured in
centimetres. Assume that $X$ follows a continuous distribution with PDF
$f(x;\boldsymbol{\theta})$ and cumulative distribution function
$F(x;\boldsymbol{\theta})$. Consider an HPS inventory compiled using a basal
area factor (BAF) $C_{BA}$. Let $\mathcal{I}$ denote the set of DBH classes and
$x_i$ the centre of class $i \in \mathcal{I}$. For an inventory comprising
plots $\mathcal{J}$, the HPS tally for class $i$ in plot $j$ is $t_{ij}$ and
the corresponding mean basal area is $\bar{g}_{ij}$. The expanded stand table
estimate of stems per hectare in class $i$ is
\begin{equation}
  \hat{y}_i = C_{BA} \frac{1}{|\mathcal{J}|} \sum_{j \in \mathcal{J}} t_{ij} \bar{g}_{ij}^{-1}.
  \label{eq:stand-table}
\end{equation}
Because the inclusion probability of a stem is proportional to its basal area,
the HPS tally is size-biased of order two. The HPS expansion factor that maps a
tally back to expected stems per hectare at DBH $x$ is
\begin{equation}
  f_E(x; C_{BA}) = \frac{40000 \, C_{BA}}{\pi x^2},
  \label{eq:expansion}
\end{equation}
with multiplicative inverse (compression factor):
\begin{equation}
  f_C(x; C_{BA}) = \frac{\pi x^2}{40000 \, C_{BA}} = f_E(x; C_{BA})^{-1}.
  \label{eq:compression}
\end{equation}

\subsection*{Reference estimator}
The reference approach \citep{van1986, ducey2015} fits a size-biased PDF
$f_{\text{sb}}(x; \boldsymbol{\theta})$ to the raw HPS tallies. For any assumed
standard-form PDF $f(x; \boldsymbol{\theta})$, the size-biased form of order
two is
\begin{equation}
  f_{\text{sb}}(x; \boldsymbol{\theta}) =
  \frac{x^2 f(x; \boldsymbol{\theta})}{\int_0^\infty x^2 f(x; \boldsymbol{\theta}) \, \mathrm{d}x}.
  \label{eq:size-biased}
\end{equation}
The typical objective minimises the sum of squared deviations between expected
and observed tallies:
\begin{equation}
  Z_{\text{C}}(\boldsymbol{\theta}) =
  \sum_{i \in \mathcal{I}} \left[t_i - s \, f_{\text{sb}}(x_i; \boldsymbol{\theta})\right]^2,
  \label{eq:control-obj}
\end{equation}
where $t_i = \sum_{j \in \mathcal{J}} t_{ij}$ and $s$ is a scaling parameter
that reconciles the continuous PDF with discrete bin counts.

\subsection*{Weighted estimator}
Our alternative estimator fits the standard-form PDF directly to the expanded
stand table while embedding the size-bias correction through bin-wise weights.
Let $\hat{y}_i = f_E(x_i; C_{BA}) t_i$ denote the expanded tally in class $i$.
The weighted objective is
\begin{equation}
  Z_{\text{T}}(\boldsymbol{\theta}) =
  \sum_{i \in \mathcal{I}} w_i^2
  \left[\hat{y}_i - s \, f(x_i; \boldsymbol{\theta})\right]^2,
  \qquad
  w_i = f_C(x_i; C_{BA}),
  \label{eq:test-obj}
\end{equation}
which mirrors a non-linear least squares problem with heteroscedastic error
variance proportional to the expansion factor.

\subsection{Equivalence}
Substituting $\hat{y}_i = f_E(x_i; C_{BA}) t_i$ and $w_i = f_C(x_i; C_{BA})$
into \eqref{eq:test-obj} yields
\begin{align}
  Z_{\text{T}}(\boldsymbol{\theta})
    &= \sum_{i \in \mathcal{I}}
       \left[f_C(x_i; C_{BA}) f_E(x_i; C_{BA}) t_i - f_C(x_i; C_{BA}) s f(x_i; \boldsymbol{\theta})\right]^2 \nonumber \\
    &= \sum_{i \in \mathcal{I}}
       \left[t_i - s \, \frac{f(x_i; \boldsymbol{\theta})}{f_E(x_i; C_{BA})^{-1}}\right]^2 \nonumber \\
    &= \sum_{i \in \mathcal{I}}
       \left[t_i - s \, \tilde{c} \, x_i^2 f(x_i; \boldsymbol{\theta})\right]^2,
\end{align}
where $\tilde{c} = \pi/(40000 C_{BA})$ is a constant that does not depend on
$\boldsymbol{\theta}$. The normalising constant in \eqref{eq:size-biased} is
also independent of $\boldsymbol{\theta}$; therefore the minimisers of
$Z_{\text{C}}(\boldsymbol{\theta})$ and $Z_{\text{T}}(\boldsymbol{\theta})$ are
the same. A full derivation is provided in the project repository and extends
to unbinned likelihood formulations.

\subsection{Computational experiment}

We evaluate the method using permanent sample plot (PSP) data from
Quebec \citep{gouvernement2019}, which provide stem tallies for 30
species--cover combinations. HPS tallies are emulated by scaling each bin by the reciprocal
expansion factor, and for three representative meta-plots (spruce-pine-fir-larch
(SPFL) softwood, SPFL-S; birch mixedwood, Birch-M; and maple hardwood, Maple-H)
we fit Weibull and Gamma distributions with both estimators. Analyses assume
fixed binning, independent class tallies, and parametric Weibull or Gamma
families for the experiment. The reproducible pipeline in the project repository
handles preprocessing and fitting.

\section{Results}

Figure~\ref{fig:comparison} compares fitted distributions across the six
meta-plot--distribution combinations. In both HPS tally space and stand-table
space the curves produced by the reference and weighted estimators are visually
similar. Table~\ref{tab:comparison} reports relative $\ell_2$ norms
for the fitted curves: in stand-table space the weighted estimator deviates by
no more than 2.4\%, and in HPS space by at most 8.5\%. Absolute residual sums of
squares (RSS) are large because they are reported on the original scale of the data,
but the percentage differences indicate that the weighting scheme reproduces the
size-biased estimator within numerical tolerance. Chi-square goodness-of-fit
statistics for both estimators remain within the same order of magnitude across
all species and cover type combinations.

\begin{figure}[htbp]
  \centering
  \includegraphics[width=\textwidth]{Fig1.pdf}
  \caption{Example comparison of control (size-biased) and test (weighted)
  estimators for the SPFL-S meta-plot (spruce-pine-fir-larch softwood). Solid lines denote weighted fits and
  dashed lines denote size-biased fits. Points represent empirical data
  (expanded stand tables on the right panels). Remaining combinations are
  available in the project repository}
  \label{fig:comparison}
\end{figure}

Although absolute RSS values are large on the native scale, relative
$\ell_2$ errors stay below 8.5\% and parameter estimates match to four
decimal places. A Jupyter notebook in the project repository documents
the full set of replicates and provides interactive diagnostics.

\begin{table}[htbp]
  \centering
  \resizebox{\textwidth}{!}{%
    \begin{tabular}{lllrrrrrr}
    \toprule
    species\_group & cover\_type & distribution & sample\_size & rel\_l2\_stand(\%) & rel\_l2\_hps(\%) & rss\_control\_hps & rss\_test\_stand & rss\_diff\_stand \\
    \midrule
    sepm & r & weibull & 152875 & 1.32 & 4.41 & 2.344e+06 & 5.174e+10 & 2.084e+10 \\
    sepm & r & gamma & 152875 & 1.51 & 4.72 & 2.432e+06 & 5.421e+10 & 2.758e+10 \\
    bop & m & weibull & 40125 & 2.31 & 8.49 & 5.866e+05 & 3.967e+09 & 2.322e+09 \\
    bop & m & gamma & 40125 & 2.31 & 7.97 & 6.028e+05 & 3.915e+09 & 2.327e+09 \\
    ers & f & weibull & 32250 & 1.90 & 4.19 & 2.369e+05 & 1.215e+09 & 5.313e+08 \\
    ers & f & gamma & 32250 & 1.86 & 7.05 & 2.163e+05 & 1.596e+09 & 5.061e+08 \\
    \bottomrule
    \end{tabular}
  }
  \caption{Residual diagnostics for control (size-biased) and test (weighted)
  estimators across species groups, cover types, and assumed distributions.
  Values shown are residual sum of squares (RSS) in native scale and chi-square statistics.
  Species-group codes are sepm (SPFL: spruce-pine-fir-larch), bop (birch), and ers (maple);
  cover-type codes r, m, f denote softwood, mixedwood, and hardwood classes.}
  \label{tab:comparison}
\end{table}

\section{Discussion}

The weighted estimator eliminates the need to derive bespoke size-biased
PDFs because it operates on the standard functions distributed with
statistical libraries. This avoids algebraic effort and restrictive
parameter bounds while preserving the statistical correctness of the
reference estimator.

Across all species and cover combinations, the weighted estimator stays
within 2.4\% of the size-biased control in stand-table space and within
8.5\% in HPS tally space (Table~\ref{tab:comparison}), so the weighted
estimator can replace the reference estimator without perceptible loss.

Appendix~A (submitted separately) derives equivalence for non-linear
least squares, and the same argument extends to likelihood-based models,
including Poisson Horvitz--Thompson weighting and potential Bayesian
variants.

The reproducible pipeline in the project repository handles data preparation,
fitting, and figure/table generation, so new datasets can be analysed by
refreshing the binned meta-plot file.

Our PSP-derived pseudo-HPS dataset assumes that expanded tables share the
variance structure of true HPS tallies; validating the method on genuine HPS
inventories is a natural next step.

\section{Conclusion}

We presented a weighted fitting procedure that reproduces the behaviour of the
traditional size-biased estimator for deriving stem diameter distributions from
HPS data. The method is easier to implement, numerically stable, and accompanied
by a fully reproducible research package. The experiments show that the
weighted and size-biased estimators yield closely aligned fits across a range
of species groups, cover types, and target distributions. The formal derivation in
Appendix~A (submitted separately) codifies this equivalence and supports
adoption of the weighted formulation in operational workflows. We recommend
using the provided templates as a foundation for future methodological
extensions.

\section*{Acknowledgements}
\noindent\textbf{Funding.} This research was supported by the Forest Resources Environmental Services Hub (FRESH) at the University of British Columbia. Funding was provided by the Natural Sciences and Engineering Research Council of Canada (NSERC, grant RGPIN-2023-04197).\\
\noindent\textbf{Competing interests.} The author declares no competing interests.\\
\noindent\textbf{Author contributions.} G.E. Paradis designed the study, performed the analysis, and prepared the manuscript.\\
\noindent\textbf{Employment.} G.E. Paradis is employed as an Assistant Professor at the University of British Columbia.\\
\noindent\textbf{Data availability.} Processed datasets and reproduction scripts are available from the project repository (\href{https://github.com/UBC-FRESH/dbhdistfit-papers}{UBC-FRESH/dbhdistfit-papers}), along with extensive metadata and notes that expand on this brief report. Raw PSP data are distributed by Donn\'ees Qu\'ebec and may be downloaded from the agency's public portal.\\
\noindent\textbf{Consent/Ethics.} Not applicable; the study uses publicly available data.\\

\bibliographystyle{plainnat}
\bibliography{references}

\end{document}
